% Generated draft: requires mathematical and source review.
\documentclass[11pt]{article}
\usepackage[margin=25mm]{geometry}
\usepackage{fontspec}
\setmainfont{Latin Modern Roman}
\usepackage{amsmath,amssymb,mathtools,unicode-math}
\setmathfont{Latin Modern Math}
\usepackage{enumitem,needspace,xurl,hyperref}
\setlength{\parindent}{0pt}
\setlength{\parskip}{5pt}
\setlength{\emergencystretch}{3em}
\newcommand{\N}{\mathbb{N}}
\newcommand{\Z}{\mathbb{Z}}
\newcommand{\Q}{\mathbb{Q}}
\newcommand{\R}{\mathbb{R}}
\newcommand{\C}{\mathbb{C}}
\newcommand{\F}{\mathbb{F}}
\newcommand{\A}{\mathbb{A}}
\newcommand{\PP}{\mathbb{P}}
\newcommand{\E}{\mathbb{E}}
\newcommand{\eps}{\varepsilon}
\newcommand{\dd}{\,\mathrm{d}}
\begin{document}
\section*{UnsolvedMath ID: 30003328}
\textbf{Source problem code:} catalogue entry
\subsection{Definitions and mathematical statement}
Let \(\mathrm{No}\) denote the ordered field of surreal numbers, whose underlying collection is a proper class. For an ordered field \(K\), a semialgebraic subset of \(K^n\) is a finite Boolean combination of polynomial equalities and inequalities with coefficients in \(K\); a function is semialgebraic if its graph is semialgebraic.

For terminology only, in the usual real setting a restricted analytic function is a function \(f\colon\R^n\to\R\) obtained from a real analytic function \(g\) defined on an open neighbourhood of \([-1,1]^n\) by setting
\[
f(x)=\begin{cases}g(x),&x\in[-1,1]^n,\\0,&x\notin[-1,1]^n.\end{cases}
\]
The structure \(\R_{\mathrm{an}}\) is the real ordered field expanded by all such functions. Globally subanalytic sets are the sets definable in this structure, and globally subanalytic functions are those with globally subanalytic graphs. These real definitions do not specify interpretations of the functions on surreal arguments.

\textbf{Source research question.} Can the surreal-number measure and integration theory referred to in the source be extended beyond semialgebraic and globally subanalytic sets and functions, in particular to restricted analytic functions and exponentiation?

\textbf{Unresolved formulation.} This is an extension programme rather than a fully specified existence proposition. The supplied record does not define the starting measure and integral, their domains and value spaces, the axioms an extension must satisfy, or the required agreement with the starting theory. It also does not specify which surreal interpretations of restricted analytic functions are intended, which exponential function is intended, or which sets, functions, dimensions, and integration bounds must be covered. Consequently, no precise universal quantifiers over these objects, or definitive existence assertion for an extension, can be extracted without additional source information. Under the conventional real definitions above, restricted analytic functions are already globally subanalytic; the intended additional surreal coverage therefore needs clarification.
\subsection{Short English statement}
The question asks for a broader measure and integration theory on surreal numbers, especially involving restricted analytic functions and exponentiation. The source leaves the starting theory and the requirements for an extension unspecified. Its supplied status is partially solved; this is a dated curation label, not an independently verified mathematical conclusion.
\subsection{Sources}
\begin{itemize}
\item Supplied local UnsolvedMath registry record, permanent ID \(30003328\), title \emph{Surreal Integration Beyond Globally Subanalytic Functions}, source problem code \texttt{catalogue entry}, in the Oberwolfach Reports Open Problems collection. The registry labels the statement \texttt{exact} and records that it was retained byte-for-byte after an internal readability check. No individual proposer is supplied.
\item Supplied source citation: \emph{Mini-Workshop: Surreal Numbers, Surreal Analysis, Hahn Fields and Derivations} (2017), Oberwolfach Reports, pp.\ 3313--3372, DOI \texttt{10.4171/owr/2016/60}. Original OWR report: \url{https://doi.org/10.4171/owr/2016/60}. The record separately gives the proposed year as 2016. Its background describes Oberwolfach Reports as documenting workshops and their open-problem discussions.
\item Supplied literature reference: Costin--Ehrlich--Friedman, \emph{Integration on the Surreals}, \url{https://arxiv.org/abs/1505.02478}.
\item Supplied literature reference: Costin--Ehrlich, \emph{Integration on the Surreals}, \url{https://doi.org/10.1016/j.aim.2024.109823}.
\item Literature assessment checked at \texttt{2026-08-21}: the supplied curation labels the problem partially solved and reports extensions from the stated semialgebraic and globally subanalytic setting to substantial analytic, meromorphic, Borel-summable, resurgent, and differential-equation classes, while describing the requested general restricted-analytic and exponential theory as remaining open. This is supplied dated curation derived from the OWR record, not an independently verified assessment. The record requests independent verification of time-sensitive status and cited claims.
\end{itemize}
\end{document}
